\section{Naming the curves}

It is useful to give a name to the curves. Let us take as an example the function $f$ defined by:

\[
   x>0\ \text{and}\ f(x)=\dfrac{x^2+x-1}{x^3}
\]

A title is easily written with the macro \tkzcname{tkzFnText}. The coordinates are given in the units of the axes. To place a text along the curve, the simplest way is to define a point of the curve and to use its name.

\begin{tkzltxexample}[num]
  \tkzFnDefPoint(f,3){P}
  \tkzFnText[above right](P){${\mathcal{C}}_f$}
\end{tkzltxexample}

The first line defines the point $P$ of the graph with $x=3$. The second one places the text at $P$. The \TIKZ\ options refine the result.

\subsection{Adding a label}

\begin{center}
\begin{tkzexample}[vbox]
\begin{tikzpicture}
  \tkzFnInit[xmin=0,xmax=10,
             ymin=0,ymax=1.2,ystep=0.2]
  \tkzFnGrid
  \tkzFnAxes
  \tkzFnDraw[name=f,thick,color=red,domain=0.55:10]{(x*x+x-1)/x^3}
  \tkzFnText(3,-0.3){\textbf{Graph of} $\mathbf{f}$}
  \tkzFnDefPoint(f,3){P}
  \tkzFnText[above right,text=red](P){${\mathcal{C}}_f$}
\end{tikzpicture}
\end{tkzexample}
\end{center}

\endinput
